% 第 54 章　原子怎样组成分子与固体
\chapter{原子怎样组成分子与固体}

第~53~章结束时，我们把一百多种元素各自的脾气都归因于一件事：原子最外层住着几个电子。可是单质原子只是世界的零件箱。你周围的几乎一切——水、盐、塑料、钢铁、玻璃、肌肉、芯片——都是原子组合而成的分子与固体。组合的方式不同，同样一堆原子就能变成天差地别的东西。本章要回答的问题只有一个：两个原子靠近时，究竟发生了什么，使它们有时紧紧抱成一团，有时漠然擦肩，有时拉帮结派地组成能导电、能折断、能拉丝的固体？答案会一路把我们带到你手机里那块芯片的心脏。

\step{反常识开场}

铅笔芯和钻戒，是同一个东西。

这句话字面上是错的——铅笔芯里混着黏土，钻石常含杂质——但把杂质剔掉之后，两者的本体是同一种元素：纯碳。金刚石里每个原子都是碳，石墨里每个原子也都是碳，连一个电子的数目都不差。可是看看这对“亲兄弟”的表现：金刚石是自然界最硬的东西，透明，不导电；石墨软得能在纸上蹭下碎屑（你写字靠的就是它），乌黑，而且导电性好到能当电极。同一种原子，一模一样的核外电子，凭什么一个硬到能切玻璃，一个软到能做润滑剂？一个对电流关门，一个对电流放行？

如果你觉得这只是巧合，再看一对。食盐和铁勺都是固体，都由原子规则地堆成晶体。可是把盐粒往桌上一磕，它应声碎裂；把铁丝弯过来、压过去、拉成细丝，它变形却不肯断。同样是“原子排好队的固体”，一个脆得像饼干，一个韧得像橡皮泥。还有更奇怪的：盐粒本身是绝缘体，泡进水里却变成导电的盐水；而铜无论熔化还是固态都是良导体。

最后一件最反常的事，留给你天天带在身上的手机。一块指甲盖大小的硅片上刻着上百亿个开关，每个开关只是同一片硅里两个挨在一起、成分略有差异的区域。没有活动的零件，没有真空管，一片看起来死气沉沉的灰色晶体，凭什么能判断“该让电流过”还是“不该让它过”，还能把微弱的信号放大成能推动耳机的声音？

\step{先猜一猜}

照例，先写下你的三个预测，再往下读。

第一猜：两个氢原子靠近会黏成氢分子，两个氦原子靠近却永远不黏。氢和氦都是电中性的原子，你觉得“黏不黏”取决于（a）原子的质量，（b）原子核的电荷数，（c）核外电子的某种排布方式？如果你选了（c），请具体猜一猜：氢有一个电子，氦有两个，差别是“一个”还是“单双”？

第二猜：碳原子有四个最外层电子。金刚石中每个碳和四个邻居相连，石墨中每个碳和三个邻居相连。仅由这两个数字，你能猜出哪一种结构里电子被“拴得更死”、哪一种有电子能到处跑吗？哪一种导电？

第三猜：把一亿亿个铜原子堆成一块铜，电流可以在里面流动。电流是电荷的移动——那么这些跑动的是哪种电荷？是铜原子核在跑，是铜离子在跑，还是有一部分电子不再属于任何一个原子、成了整块金属的“公共财产”？如果是最后一种，它们为什么不会一股脑从金属里漏出去？

\step{动手实验}

\begin{experiment}[铅笔芯是导体吗]
\textbf{材料}：一支铅笔（2B 以上更软、石墨更多，效果更好），一节 9~V 电池（或三节 1.5~V 电池串联），一只发光二极管（LED，电子市场几分钱一只），一小撮食盐，一点白糖，两杯水，几根导线（没有导线可用铝箔搓成条）。

\textbf{第一步}：把铅笔两头削尖，露出笔芯。用导线把电池、LED 和铅笔芯串成一个圈——电流必须从电池出发，穿过笔芯和 LED，再回到电池。LED 有方向性，不亮就把它的两条腿对调。你会发现 LED 亮了：铅笔芯是导体，电流从石墨里穿了过去。把笔芯换成橡皮、塑料尺，灯立刻熄灭。

\textbf{第二步}：把两根导线的一端同时插进一小堆干燥的食盐里（彼此别碰上），另一端照旧接电池和 LED。灯不亮——固态的盐是绝缘体。

\textbf{第三步}：把盐搅进一杯水里，再把两根导线插进盐水（同样别碰上）。灯亮了。换成白糖水，灯不亮。

\textbf{记录与思考}：石墨和盐都是原子整齐排列的晶体，一个导电一个不导电；同一撮盐，干燥时绝缘、溶进水里导电；糖水和盐水看起来一样清，导电性却天差地别。导电的到底是“固体”还是“液体”？是“原子”还是某种从原子里解放出来的东西？请写下你的解释，并回答：盐水中跑动的电荷，和铅笔芯中跑动的电荷，是同一种东西吗？
\end{experiment}

这个实验里藏着本章的全部悬念。提示一句：盐水里跑动的是钠离子和氯离子——带电的整个原子；而铅笔芯里跑动的只是电子，原子本身待在原地一动不动。同样是“导电”，微观图景完全不同，固体导电的那条路才是本章的主角。

\step{旧模型遇到麻烦}

先用我们手上的旧装备——库仑定律加上第~52~章的原子图像——认真攻打这些问题，看在哪里翻车。

第一场翻车：静电说不清“谁跟谁黏”。氢分子由两个氢原子组成，每个氢原子整体电中性。两个中性的东西，库仑力应当相互抵消，至多剩下电子云极化带来的微弱吸引——那点吸引弱到连氦都冻不成液体，绝撑不起一根拉不断的化学键。更打脸的是选择性：氢和氢黏，氢和氦不黏，氦和氦不黏。静电力只认电荷，中性原子在它面前长得一模一样，它没有任何机制区分“该黏”和“不该黏”。经典物理甚至给不出“分子”这个概念存在的理由——原子为什么不永远各过各的？

第二场翻车：就算承认原子会黏，旧模型说不清“黏几个”和“往哪黏”。一个氢原子只黏一个氢原子，绝不黏第二个；一个氧原子偏要黏两个氢，而且三根键之间夹角固定在一百零四度半。如果成键只是静电吸引，为什么吸引会“吃饱”——饱和之后再来一个原子就坚决不要？为什么吸引还带方向？一根无方向的吸引力应该来者不拒、四面八方均匀才是。旧模型对此只能支支吾吾。

第三场翻车：固体性质的巨大分裂毫无头绪。把原子想象成小硬球堆起来，再假设球与球之间有些说不清来源的“黏性”——这就是旧模型对固体的全部理解。可这个模型预言所有晶体应当大同小异，现实却是：银的导电能力大约是金刚石的一万亿亿倍（电阻率相差二十个数量级以上）；石墨沿层面方向的导电能力是垂直层面方向的数百倍；金刚石硬得天下第一，石墨却是一等一的固体润滑剂。一堆“黏起来的球”产生不了二十个数量级的分裂，更解释不了为什么导电与否偏偏和硬度、透明度捆绑出现——金刚石又硬又透明又不导电，这三件事之间有什么共谋？

三场翻车指向同一个缺口：我们一直在用经典的眼光看电子，把电子当成待在原子里的小颗粒。可第~52~章已经说过，电子在原子里的居所是由量子数编址的量子态，是一片有形状、有能量的电子云。两个原子靠近时，真正发生的故事是这些量子态的重新洗牌。

\step{发明新概念}

现在像物理学家一样，为填上缺口发明四样新东西。

第一样：\textbf{分子轨道}。两个原子靠近到电子云开始重叠时，电子面对的局面变了：附近有两个原子核，都是正电荷的“深谷”。量子力学的账本不允许原来的两套原子轨道原封不动地并存，它们必须重新组合成属于整个分子的新量子态——分子轨道。关键事实是这样的：两个原子轨道组合，总是给出两个新轨道，一个能量比原来低，叫\textbf{成键轨道}，一个能量比原来高，叫\textbf{反键轨道}。成键轨道的电子云堆在两个核之间，负电荷像一座桥横在峡谷上，同时被两边的正核拉住；反键轨道的电子云在两核之间恰恰有一个空档，电子待在那里什么也黏不住，还推高能量。这像什么？像两盏靠近的灯，灯光叠加处更亮、交界处也可能因波动叠加出现暗纹——电子云是概率的波，两个波叠在一起，有的地方概率涨起来，有的地方退干净。它不像什么？它不像两颗球用弹簧连起来——分子轨道不是两个核之间的一根实物纽带，它是电子在整个分子范围内的一个量子态，电子不“沿着它跑”，而是“以它为分布”。

有了分子轨道，氢与氦的谜题当场解开，而且答案精确地就是你在“先猜一猜”里被问到的那一点——单双。氢分子有两个电子，按第~53~章的泡利不相容原理，自旋相反的两个电子可以住进同一个成键轨道：两个电子都住进了“降价房”，总能量低于两个孤零零的原子，分子稳定。氦原子各带两个电子，氦二分子要安置四个电子：成键轨道住满两个之后，剩下两个被迫住进反键轨道。两个电子省下的能量被另两个电子花光的能量抵消（实际还略亏），总能量不降反升——于是氦分子不存在。成键不是“原子喜欢在一起”，而是量子态重排之后有没有空着的低价房，以及住进去划不划算。泡利原理还顺手解释了饱和性：氢分子的成键轨道只有两个座位，第三个氢原子再来，只能住反键轨道，不请自来还要倒贴能量，于是被拒之门外。化学键的“吃饱”，是座位表说了算。至于方向性——水分子的夹角——来自原子轨道本身的形状：第~52~章的 p 轨道是有朝向的哑铃，成键时哑铃对准谁，键就指向谁。

第二样：\textbf{离子键}——电子干脆整体搬家。钠最外层孤零零一个电子，拽走它只要约 \(5.1\)~eV（电离能）；氯最外层只差一个电子就满员，塞给它一个电子反而倒找约 \(3.6\)~eV（电子亲和能）。账面上看，把一个电子从钠搬到氯要花 \(1.5\)~eV，似乎不划算——单看这一步，直觉会说反应根本不会发生。真正的收益在下一步：搬完家，钠成正离子、氯成负离子，异性相吸，而且不是一对两对地吸——在食盐晶体里，每个离子被六个异号离子贴身围住，静电吸引的账要按整块晶体来算。这份\textbf{晶格能}摊到每对离子约 \(8\)~eV，减去搬家成本，净赚超过 \(6\)~eV。离子键不是“一个钠拉住一个氯”，而是全体离子共同买下的静电团购。这也立刻解释了盐的脆：晶体里同号离子与异号离子交替排列，像一格白一格黑的棋盘；用力推晶体让一层滑过半格，同号离子骤然面对面，吸引变排斥，晶体沿这一面干脆地断开——盐粒应声碎裂。它还解释了实验第三步：固态盐里离子被晶格锁在原地，无法移动，所以不导电；溶进水里，离子脱笼而出，整杯水里都是能跑的电荷，灯就亮了。

第三样：\textbf{金属键}——离域的极致。金属原子的最外层电子拴得不紧，当成亿个金属原子堆在一起时，最外层的原子轨道层层叠叠地重组，结果是：这批电子的活动范围不再是某个原子，也不是某对原子，而是整块金属。它们成了一片在正离子晶格的骨架间流动的“电子海”。这像什么？像一整座城市的地铁系统：任何乘客都不固定属于哪一站，全线网通行。它不像什么？不像自由气体——电子不是漫无目的地乱撞一气后就散掉，它们被所有正离子共同的静电引力留在金属内部（这正是第三猜的答案：电子漏不出去，因为走出表面等于离开全体正离子布下的引力盆地，要付一笔能量，这笔账叫逸出功）。金属的招牌性质全部从这片海读出：电子海无方向地托住正离子骨架，所以原子面之间滑移时键不挑方向、不断裂——金属能弯曲、能拉丝，这就是铁勺弯而不碎的原因，也是它与食盐脆性的根本差别；电子海里有大量几乎自由的载流子，电场一来就能定向漂移——金属导电；电子跑得勤快，把热量从一头搬到另一头——金属导热。

第四样：\textbf{能带}——从两个原子到一亿亿个原子。两个原子靠近时一个能级劈成两个（成键与反键），三个原子就是三个，\(N\) 个原子堆成固体，每个原来的原子能级就劈成 \(N\) 个间距极小的能级，密密麻麻连成一条能量上连续的带子。一块一立方厘米的铜大约有 \(10^{23}\) 个原子，这些能级挤在宽度只有几个 eV 的范围里，相邻间距小到 \(10^{-22}\)~eV 量级——任何仪器都分不出缝隙，于是固体里的电子能级不再是离散的台阶，而是一条条连续的\textbf{能带}，带与带之间可能隔着电子绝对无法落脚的禁区，叫\textbf{带隙}。能带是理解一切固体电性的钥匙，它的详细账目放到下一节细算。

在收进新工具之前，必须郑重声明一件事：化学键不是一种新发现的力。从头到尾，舞台上只有电磁相互作用这一个演员，但剧本是量子力学的——全同粒子的记账规则、泡利不相容原理、电子云的离域与重组。键是电磁力按量子规则演出的集体节目。这个声明不是客套：它意味着解释分子和固体不需要任何新物理，也意味着这一切原则上可以计算——尽管实际的计算难到养活了整个计算化学行业。

\step{一句话模型}

\begin{oneline}
原子结合成分子与固体，本质是电子量子态的重新洗牌：两个原子的轨道组合出一对降价房（成键轨道）与高价房（反键轨道），电子按泡利原理从低价房住起，住得划算就成键；一亿亿个原子一起洗牌时，能级密集成能带——电子填满的带是死的，填了一半的带里电子能自由流动，带隙宽窄决定这块固体是导体、半导体还是绝缘体。
\end{oneline}

这个模型像什么？像剧院的售票逻辑：座位分楼层（能带），同一楼层票价相同，一人一座（泡利原理），从便宜票卖起；坐不满的楼层里观众可以换座走动（导电），坐满的楼层纹丝不动，想动只能买更高楼层的票（越过带隙）。它不像什么？第一，电子不是真的观众——它们不挑座位，是量子规则规定了一切占据方式，模型说的是“哪些态被占”，不是“电子的选择”。第二，“走动”不是宏观意义上的走动：满带里电子并非静止，它们也在运动，只是所有运动恰好互相抵消，不产生净电流——导电的判据是“有没有相邻的空态可去”，不是“电子动不动”。第三，票价图（能带结构）是整块晶体共同的性质，不是每个原子各有一张——把固体想成一堆独立小原子各自为政，正是旧模型翻车的根源。

\step{数学翻译器}

先做第一件翻译：把“成键划算吗”画成一条曲线。设两个原子核之间的距离为 \(r\)，体系总能量为 \(U(r)\)。\(r\) 很大时两个原子互不相干，记 \(U=0\)；\(r\) 缩小到电子云开始重叠，成键轨道被占据，能量下降；\(r\) 继续缩小，两核之间的裸排斥、以及泡利原理逼电子上高价房的代价占了上风，能量掉头向上、急剧飙升。中间存在一个能量最低的距离 \(r_0\)——这就是键长；从谷底到零的深度 \(D\)——这就是键能。

\begin{figure}[htbp]
\centering
\begin{tikzpicture}[scale=1.0,>=Stealth]
  \draw[->] (0,0) -- (8.2,0) node[below left] {核间距 \(r\)};
  \draw[->] (0,-3.2) -- (0,2.0) node[above left] {能量 \(U\)};
  \draw[thick,blue!60!black] (0.45,1.9)
    .. controls (0.75,-1.4) and (1.1,-2.55) .. (1.7,-2.62)
    .. controls (2.6,-2.7) and (3.6,-1.5) .. (5.2,-0.45)
    .. controls (6.4,-0.15) .. (8.0,-0.02);
  \draw[dashed] (1.7,-2.62) -- (1.7,0) node[above] {\(r_0\)};
  \draw[dashed] (0,-2.62) -- (1.7,-2.62);
  \draw[<->] (-0.12,-2.62) -- (-0.12,0) node[midway,left] {\(D\)};
  \node[blue!60!black] at (5.6,-1.6) {成键分子的势能曲线};
  \draw[thick,red!60!black] (0.55,1.9)
    .. controls (1.2,0.6) and (2.2,0.25) .. (4.0,0.12)
    .. controls (6.0,0.04) .. (8.0,0.01);
  \node[red!60!black] at (4.6,1.1) {不成键（如 He--He）};
\end{tikzpicture}
\end{figure}

用特殊情况检查这条曲线，每一个都该符合你的直觉：\(r\to\infty\) 时 \(U\to 0\)，距离无穷远，互不相干，对；\(r\to 0\) 时 \(U\to+\infty\)，两个带正电的核几乎贴脸，排斥无穷大，对；\(r=r_0\) 处曲线斜率为零——力是能量曲线的坡度，谷底坡度为零，原子不受力，正好停在那里，对；\(r\) 比 \(r_0\) 略小时曲线斜率为负，力指向 \(r\) 增大的方向，把原子推回去，对；\(r\) 比 \(r_0\) 略大时斜率为正，力把原子拉回来，对。谷底就是一个稳定的平衡位置。而那条氦的曲线没有谷底，一路只是排斥加微弱的尾巴——不成键。

再看键能的真实数量级。典型的共价键键能是几个电子伏特：H--H 约 \(4.5\)~eV，C--C 约 \(3.6\)~eV；食盐里每对离子的净结合收益约 \(6\)~eV。拿它和室温的热运动比：热学篇里说过，\(300\)~K 时每个自由度摊到的热运动能量约 \(kT\approx0.026\)~eV。化学键比室温的热骚动硬一百倍以上——所以桌子、水、你的骨骼在室温下是稳定的东西，分子不会被热运动随手撞散。但反过来说，把温度抬高到几千开（\(kT\) 涨到零点几~eV），再算上高能尾巴里少数格外活跃的分子，化学键就开始大批断裂——这就是燃烧和高温冶炼能改写物质的能量根源。一个电子伏特与零点零二六个电子伏特的比值，悄无声息地划出了“化学世界稳定”与“物质瓦解”的温度边界。

第二件翻译：一粒盐里有多少个离子，排起来有多大？取一粒边长约 \(0.5\)~mm 的食盐，质量约 \(1\)~mg。氯化钠的摩尔质量约 \(58.5\)~g/mol，所以这一粒盐约含 \(10^{-3}/58.5\approx1.7\times10^{-5}\)~mol 的离子对，乘上阿伏伽德罗常数，约 \(10^{19}\) 对离子。晶体中相邻钠离子与氯离子的间距是 \(0.28\)~nm，把 \(10^{19}\) 对离子堆成立方体，每边约 \((10^{19})^{1/3}\approx2\times10^{6}\) 个，乘上间距，约 \(0.6\)~mm——正好就是那粒盐的尺寸。账算平了：一粒看得见的盐，每个方向上都排着两百万个离子。这就是“宏观一小粒、微观天文数”的直接换算，也是为什么单个化学键只有几个电子伏特、掰碎一块晶体却要宏观可观的力气——你一次折断的是亿亿根键的团购。

第三件翻译：能带有多“连续”？一立方厘米的铜约含 \(8.5\times10^{22}\) 个原子，最外层那个能级劈成 \(8.5\times10^{22}\) 个子能级，全部挤在宽度约几个电子伏特的带内。相邻子能级的间距约 \(10^{-22}\)~eV——是室温 \(kT\) 的一百亿亿分之一。别说仪器，连最微弱的热扰动都足以把电子从一个子能级扫到相邻子能级。所以“带内连续”不是近似口号，而是任何物理过程都无法分辨其缝隙的事实。而真正分立的是带与带之间：带隙里没有可供电子落脚的态，一个态都没有。

\begin{figure}[htbp]
\centering
\begin{tikzpicture}[scale=1.0,>=Stealth]
  % 单原子：一条能级
  \draw[thick] (0,0) -- (1.2,0);
  \node[below] at (0.6,-0.1) {一个原子};
  \node[right] at (1.3,0) {一条能级};
  % 两个原子：劈成两条
  \draw[thick] (3.2,0.25) -- (4.4,0.25);
  \draw[thick] (3.2,-0.25) -- (4.4,-0.25);
  \node[below] at (3.8,-0.35) {两个原子};
  \node[right] at (4.5,0) {劈成两条};
  % N 个原子：连成能带
  \foreach \y in {-0.5,-0.4,-0.3,-0.2,-0.1,0,0.1,0.2,0.3,0.4,0.5}{
    \draw (6.6,\y) -- (7.8,\y);}
  \node[below] at (7.2,-0.6) {\(N\) 个原子};
  \node[right] at (7.9,0) {连成能带};
  \draw[->,thick] (1.9,0) -- (2.9,0);
  \draw[->,thick] (5.0,0) -- (6.0,0);
\end{tikzpicture}
\end{figure}

有了能带，导体、绝缘体、半导体的分界只需一句话。泡利原理要求电子从最低能级往上填。如果最上面的带恰好填了一半（铜、银正是如此），被占的态旁边就是空态，电场轻轻一推，电子就能挪动过去形成定向流——导体。如果最上面的带填得满满当当（叫满带），再往上隔着一条宽宽的带隙才是空带，电子身边无态可去，全体运动互相抵消，电流为零——绝缘体；金刚石的带隙约 \(5.5\)~eV，可见光光子最高约 \(3\)~eV，连光都无法把电子推过带隙，所以光直接穿过：金刚石透明，正是它不导电的同一枚硬币的另一面。半导体是带隙窄的绝缘体——硅的带隙约 \(1.1\)~eV，室温下少数电子被热运动抛过带隙，导电性介于两者之间，而且对温度和杂质极其敏感。最后这一条不是缺点，是礼物：往纯硅里掺入百万分之一量级的五价原子（如磷），每个磷多出一个用不掉的电子，成为现成的载流子（n 型）；掺三价原子（如硼），每个硼缺一个电子，留下一个能被邻居电子填补的空位，空位在电场下逆向移动，等效于正电荷在跑，叫空穴（p 型）。导电性的开关第一次握在了人手里。

把一个 n 型区和一个 p 型区做在同一块硅上，交界处就构成二极管。在更细的定量描述之前，先给主线直觉：接正向电压时，外电场把两边的载流子推向界面，电子与空穴在界面处相遇复合，新的载流子源源不断补进来，电流畅行；接反向电压时，外电场把两边的载流子都拉离界面，界面附近形成一段没有载流子的“无人区”，电流被掐断。一片没有活动部件的死晶体，就这样实现了单向放行。检查特殊情况：正向电压为零时无电流（无源无流，对）；反向电压极大时，强电场终会把价电子硬扯过带隙，二极管击穿——单向性不是绝对的，每种器件都有极限，对。在二极管旁再加一个控制极，用电场调节导电通道的宽窄，就得到晶体管：控制极上的微小电压变化，可以放行或截断大得多的电流——放大与开关，数字世界的两块基石。你手机里的百亿个晶体管，全是这一句话的工业复制。

\begin{mathbridge}
分子的键像弹簧，但弹簧的能量是量子化的。把势能曲线在谷底附近近似成抛物线（这正是振动章节里简谐近似的做法），分子振动的能级为
\[
E_n=\Bigl(n+\tfrac12\Bigr)\hbar\omega,\qquad n=0,1,2,\dots
\]
其中 \(\omega=\sqrt{k/\mu}\)，\(k\) 是键的“倔强系数”，\(\mu\) 是两个原子的约化质量。检查特殊情况：\(n=0\) 时 \(E_0=\hbar\omega/2\neq0\)——零点能，分子即使在绝对零度也不可能完全静止，否则位置和动量将同时确定，与不确定性原理冲突（这是量子账本的硬约束，与仪器好坏无关）；\(n\) 很大时相邻能级间距 \(\hbar\omega\) 相对 \(E_n\) 趋于零，量子台阶密得看不出阶梯，回到经典的连续弹簧，对。典型分子的 \(\hbar\omega\) 约零点几~eV，对应红外光子的能量——所以红外线照过来，分子能被推上一级振动台阶，这正是红外光谱识别分子的原理，也是温室气体吸收地面热辐射的入口。

分子还会转。双原子分子绕质心转动，转动惯量为 \(I=\mu r_0^2\)，转动能级为
\[
E_\ell=\frac{\hbar^2}{2I}\,\ell(\ell+1),\qquad \ell=0,1,2,\dots
\]
检查特殊情况：\(\ell=0\) 时 \(E_0=0\)，转动可以真正静止（与振动不同——转动静止并不同时固定一对共轭变量，不与不确定性原理冲突）；相邻能级间距随 \(\ell\) 线性增大，\(\ell\) 越大台阶越高，这是 \(\ell(\ell+1)\) 而非 \(\ell^2\) 的细微后果，对。HCl 分子的转动台阶间距约 \(10^{-3}\)~eV 量级，对应微波光子的能量——微波炉加热食物，正是微波光子一批批被水分子的转动态吸收，再经碰撞摊平成热运动。振动台阶在红外，转动台阶在微波，两者相差两个数量级：用一束光扫过频率，读出哪些频率被吸收，就同时读出了分子的键强、键长和质量——分子光谱学就是这么一本账。
\end{mathbridge}

\begin{challenge}
为什么满带一定不导电？直觉论证“没有空态可去”之外，还有更对称的说法。晶体中电子的量子态由晶格周期势决定，定态是被晶格调制的平面波（布洛赫波），用晶体动量 \(\hbar k\) 标记；满带中所有 \(k\) 态全被占据，而 \(k\) 与 \(-k\) 总是成对出现、速度恰好相反，无论加不加电场（只要电场不足以把电子推过带隙），占据结构在 \(k\) 空间整体平移后仍是对称填满的——总电流恒等于零。导电的本质不是“电子在动”，而是“占据情况在 \(k\) 空间不对称”。这一条在本书后面还会回来：超导正是让电子配成对、以另一种方式获得不耗散的定向流动。另外提醒：本节把晶格周期势当作给定背景，单电子近似忽略了电子之间的实时关联——对铜、硅这很好用，但对某些氧化物（不少高温超导体的母体正在此列）会错得离谱，那类材料叫莫特绝缘体，按能带图它们本该是金属。能带论是模型，不是定律；它的边界就是关联电子物理的入口。
\end{challenge}

\step{模型的边界}

本章给了三种键、一个能带，它们各自的适用边界必须标清，否则下一章你就会用错地方。

分子轨道图像的边界：它把电子分配到属于整个分子的轨道上，对氢分子、氧分子、苯这类电子明显离域的体系极其出色；而对“电子主要待在两个特定原子之间”的朴素图景，它和另一种语言——价键理论——各有侧重。两者是同一个量子力学的两种近似表述，预测一致时都对，预测不一致时都要向更完整的计算低头。请把“分子轨道”记成数学模型，不要记成电子云的实拍照片。

三种键的边界：真实世界的键很少是纯种的。水分子里的氢氧键是共价键，但氧把电子云拉得偏自己一头，使分子一头偏负、一头偏正，于是不同分子的氢与氧之间还会互相轻扯——这就是氢键。氢键只有零点几~eV，约为共价键的十分之一，却决定了冰比水轻：冰里氢键把水分子撑成空旷的架子，融化时架子部分塌掉，水反而更密——这就是冰浮在水上、湖面从顶往下冻、水下生物得以过冬的原因；氢键也托住了 DNA 双螺旋的横档。比氢键更弱的还有所有中性原子、分子之间的范德瓦尔斯吸引——它正是让氦在极低温下也能液化、让壁虎能爬上玻璃的那股力。它们都不在“离子、共价、金属”三分法里；三分法是地图的主干道，不是全部街道。

能带模型的边界：它假设原子排成完美的周期阵列。玻璃、塑料这类非晶固体没有严格周期性，能带语言要打折扣使用（玻璃照样透明，因为吸收光子的能量门槛仍在）；温度升高时原子热振动打乱周期性，电子被振动散射，金属电阻随温度上升；而把某些金属冷到极低温，电子两两配对、电阻骤然归零——超导——这是能带论必须加上新的配对机制才讲得通的故事，是本书后文的题材。莫特绝缘体的警告见挑战层，不再重复。

二极管、晶体管模型的边界：主线里“推过去就通、拉离开就断”的直觉丢掉了不少细节——正向电压不到约 \(0.7\)~V（硅）时电流其实微乎其微，所谓“开通”是指数式的陡峭爬升而非闸门式的突变；反向也总有纳安量级的漏电流。工程上用得更细，但定性方向不变。还有一点常被误解：晶体管不是“用小电流驱动大电流”的简单放大器，它是用控制极的电场改变导电通道的通导能力——改变的是材料本身的状态，能量一律由电源提供，放大不违反任何能量账。

\step{回头解释开场现象}

现在回到那三对反常的兄弟，一件件结案。

金刚石与石墨。碳最外层四个电子。在金刚石里，每个碳的四个轨道指向四面体的四个顶点，与四个邻居各成一根结实的共价键，所有电子全部住进成键轨道——满员。三维的四面体网向四面八方锁死，所以金刚石最硬；电子全部住进满带，带隙 \(5.5\)~eV，无载流子，所以绝缘；可见光光子能量不够跳过带隙，所以透明。三件事是同一件事的三张脸。在石墨里，每个碳只和同一平面内的三个邻居成键，第四个电子没有固定搭档——它离域到整个平面，成为可以在层内自由流动的电子。于是层内是强共价网（单论层内，石墨其实比金刚石还结实），层内导电；而层与层之间只剩微弱的范德瓦尔斯力，一推就滑——软、能写字、能做润滑剂；层间电子难以跨越，所以导电性沿层内与垂直层内相差数百倍。那个让直觉栽跟头的反例也在这里：直觉说“导电好的导热也好”（金属确实如此），金刚石偏不——它不导电，却是室温下导热最好的材料之一，约为铜的五倍。它搬热量不靠电子，靠晶格振动在刚性的四面体网里飞速传递。导电与导热在金属里是同一个载体干的活，所以黏在一起；在金刚石里载体不同，就分道扬镳。直觉没错在金属里，错在它把金属的偶然当成了所有固体的必然。

食盐与铁勺。食盐晶体的离子键靠全体离子的静电团购，方向上精打细算：一层滑过半格，同号离子对脸，吸引翻成排斥，晶体沿面断开——脆。铁里的金属键靠无方向的电子海，原子面滑移时电子海随时重新填平缝隙，键不因滑移而翻脸——韧。于是盐粒一磕就碎，铁丝弯而不折；于是盐能溶进水里成为导电的离子溶液（晶格被水分子拆解，离子脱笼而出），而铜以电子海导电，与是否泡水无关。

手机芯片。硅的窄带隙让它对掺杂敏感；掺进去的杂质原子送去电子或留下空穴；n 型区与 p 型区的交界构成二极管，单向放行；加控制极成为晶体管，用电场开关电流。百亿个开关的通断组合成逻辑，逻辑组合成你正在读的这块屏幕背后的一切。一块死气沉沉的灰色晶体能“判断”，靠的不是机关，而是能带图里那一条一电子伏特宽的缝隙。

\step{脑洞实验与挑战题}

\begin{thought}[把压力拧到头，化学键还活着吗]
地球上的化学键是几个~eV 的生意，原子间距由电子云的平衡位置说了算。现在把压力一路拧上去：深海底部压强约 \(10^{8}\)~Pa，物质的化学面貌基本不变；地心约 \(3\times10^{11}\)~Pa，原子被压小，有些绝缘体被压成金属，但电子仍属于各原子的壳层。继续拧到白矮星内部——压强超过 \(10^{22}\)~Pa。原子核被挤到间距远小于正常原子半径，电子不再可能属于任何一个核：泡利原理仍然生效，但“原子”这个容器已经碎了，电子成为全体原子核共同背景下的一片简并电子海，靠“一个态只住一个电子”的那股斥力顶住引力的碾压。眼熟吗？这正是第~53~章撑起桌面的那堵看不见的墙，只是这次它撑起的是一整颗死掉的恒星。请想：在这个极限里，“化学键”“分子”“固体”这些概念是在哪一步死掉的？是渐变地死，还是有明确的墓碑？再往上拧，电子被压进质子变成中子（中子星），连电磁相互作用都退位了。化学是低压宇宙的特权——这句话你现在能从电子伏特与压强的账里自己读出来。
\end{thought}

\begin{puzzles}
\item 氢分子的成键轨道可以住两个电子（自旋相反），那三个氢原子能不能共用一条“成键轨道链”，结成稳定的中性 H\(_3\) 分子？请从泡利不相容原理和成键、反键轨道的数目出发推理。（思路提示：三个原子轨道组合出三个分子轨道，按能量排队，第三个电子该住哪里？它省下的能量够不够付房钱？真实世界里 H\(_3^+\) 存在而中性 H\(_3\) 不存在，你的推理能预言这一点吗？）
\item 金属导电靠电子海，可电子带负电，全体电子定向流动时，整块金属为什么不会积累起净电荷、把自己推开？（思路提示：想想正离子骨架的电荷总量，再想想导线两端接着什么；电流是“流过”而不是“堆积”——电荷在回路中连续流动，被消耗的是能量，不是电荷。）
\item 玻璃和金刚石都不导电，都透明；纯净的食盐晶体也不导电，同样透明。请用能带语言说出这三件事的共同原因，再回答：为什么同样是绝缘体的陶瓷碗通常不透明？（思路提示：透明只要求可见光光子推不动电子，与原子是否排成完美阵列无关；想想光在不均匀的界面上除了被吸收，还会遭遇什么——散射。）
\item 一块纯净硅冷却到绝对零度附近，是导体还是绝缘体？掺入百万分之一的磷之后呢？由此回答：半导体的“半”究竟指什么？（思路提示：绝对零度下没有热运动把电子送过带隙，满带全满、空带全空；那掺杂原子送出的那个多余电子待在哪个带里？“半”不是导电能力的一半，而是导电性由谁、以什么方式决定。）
\end{puzzles}

本章把散落的原子缝成了分子与固体，用的是同一套量子规则加上电磁相互作用，没有引入任何新力。下一章我们要让这台缝好的机器运转起来：光打到原子和固体上，电子在能级与能带之间被光子推着跳——受激辐射、激光、发光二极管，还有拿原子跃迁当钟摆的原子钟。你已经知道了台阶在哪里，下一步是看什么东西在台阶上跳舞。
